\documentclass[11pt]{article}
\usepackage[T1]{fontenc}
\usepackage{lmodern,amsmath,amssymb,amsthm,mathtools,booktabs}
\usepackage[margin=1.05in]{geometry}
\usepackage[colorlinks=true,linkcolor=blue!45!black,urlcolor=blue!55!black,citecolor=blue!45!black]{hyperref}
\usepackage{xcolor}
\usepackage{microtype}
\newtheorem{theorem}{Theorem}[section]
\newtheorem{lemma}[theorem]{Lemma}
\newtheorem{corollary}[theorem]{Corollary}
\newtheorem{proposition}[theorem]{Proposition}
\theoremstyle{remark}\newtheorem{remark}[theorem]{Remark}
\newcommand{\Rea}{\operatorname{Re}}
\newcommand{\dd}{\mathrm{d}}
\newcommand{\Z}{\mathbb Z}
\newcommand{\R}{\mathbb R}
\newcommand{\C}{\mathbb C}
\newcommand{\one}{\mathbf 1}
\title{Certified Zero Counts from Real Prime Data\\[3pt]\large Rational-binomial filters for Dirichlet $L$-functions after GRH}
\author{Research note prepared with ChatGPT}
\date{26 November 2026}
\begin{document}
\maketitle
\begin{abstract}
The critical-line theorem for Dirichlet $L$-functions turns the logarithmic derivative of a paired completed function, in a squared coordinate, into a positive Stieltjes transform. We prove an explicit two-sided, finite-sample bound on the number of zeros of any Dirichlet $L$-function in the open critical strip with absolute ordinate at most $T$. Only real logarithmic-derivative values to the right of $1$ are needed; no derivative of these values, analytic continuation by computation, or prior zero location is required. A rational interpolation formula evaluates binomial tail filters exactly in terms of the samples. A correction proportional to a single trace controls the infinitely many high zeros. The bounds converge to the strict and inclusive counts, respectively, and certify the exact integer count at every zero-free boundary, even if all sampling points lie in an arbitrarily short interval. A gap-dependent error estimate and reproducible interval-arithmetic demonstration are included. The method is mathematically rigorous relative to the cited GRH theorem and classical completion theory; priority for the particular certificate is not established by the limited literature search.
\end{abstract}

\section{Source analysis and scope}
The repository \cite{repo}, inspected through its GitHub interface at main commit \texttt{72fead102f66} (full hash in the source audit), is a shortened rewriting of the Alp\"oge--Claude formal proof announced on 20 November 2026. Its README reports Lean v4.35.0, Mathlib commit \texttt{56a005855ea1}, and standard axioms only. Michael Rubinstein reports independent checks of the original proof and links to checks of the shorter proof \cite{rubinstein}. This note does not reproduce those builds; human review of the entire new proof was not claimed by the repository. The linked Wikipedia article \cite{wiki} reflects the recent proof and reviews the relation of zeros to the distribution of primes.

The readable edition identifies a chain of implications:
\begin{enumerate}
\item \texttt{grh/src/New/Reduction/MainInequality.lean}: upper size and index bounds, a spanning estimate, and a lower covolume bound yield \texttt{star\_inequality}, namely $m\,\mathrm{rem}\le fa\log^5 b$ in the repository's own parameter notation. The principal term cancels in the two-sided comparison.
\item \texttt{grh/src/New/Reduction/OneSidedBound.lean}: for $q_0$ a positive multiple of $18480$, a prime $d\equiv q_0/2-1\pmod{q_0}$, and each $\epsilon>0$, \texttt{eventually\_primeRemainder\_le\_rpow} bounds the weighted prime remainder above by $y^{-1/2+\epsilon}$ eventually. Here the remainder subtracts $\log y/\varphi(q_0)$ and a constant from a sum of $\Lambda(n)(1-\chi_d(n))(1/n-1/y)$ over $n\le y$, $n\equiv1\pmod{q_0}$, with $\chi_d$ the quadratic character modulo $d$. A one-sided nonvanishing criterion then gives \texttt{New.grh}.
\item \texttt{grh/src/New/Main.lean}: functional-equation symmetry converts the half-plane theorem into the critical-line conclusion. The statements in \texttt{grh/src/Statements.lean} are given in Mathlib's terms, not in the auxiliary geometric language.
\end{enumerate}
The new argument uses the last theorem as a proved input and takes inspiration from the two-sided comparison and one-sided analytic estimates of the earlier stages. It does \emph{not} port the theta-lift or determinant machinery itself. In the reverse direction, GRH gives positivity of an explicitly computable spectral measure; two elementary pointwise comparisons certify a discrete zero count. The scope is Dirichlet $L$-functions, not general automorphic or Dedekind $L$-functions.

\section{The positive spectral representation}
For a primitive Dirichlet character $\chi$ of conductor $q>1$, write $\chi(-1)=(-1)^\kappa$ with $\kappa\in\{0,1\}$ and set
\[
 G_\chi(s)=\left(\frac q\pi\right)^{(s+\kappa)/2}
 \Gamma\left(\frac{s+\kappa}{2}\right)L(s,\chi).
\]
For the primitive character of conductor $1$, use instead
\[
 G_1(s)=\xi(s)=\tfrac12 s(s-1)\pi^{-s/2}\Gamma(s/2)\zeta(s).
\]
The standard completion theorem says these are nonzero entire functions of order at most one, their zeros (with multiplicities) are exactly the zeros of the corresponding $L$-function in the open critical strip, and
\begin{equation}\label{eq:fe}
 G_\chi(s)=w_\chi G_{\bar\chi}(1-s),\qquad
 G_{\bar\chi}(\bar s)=\overline{G_\chi(s)},\qquad
 w_\chi w_{\bar\chi}=1.
\end{equation}
The boundary assertion uses classical nonvanishing at $\Rea s=1$, the Euler product for $\Rea s>1$, and the functional equation; the pole-cleared $\xi$ is nonzero at $0,1$. These are standard facts of Dirichlet $L$-function theory (see \cite{davenport,dlmf}). In particular, the gamma poles do not produce extra zeros of $G_\chi$. If $\chi$ is imprimitive, pass to its inducing primitive character $\chi^*$: the extra Euler factors $1-\chi^*(p)p^{-s}$ are nonzero for $\Rea s>0$, so the strip zeros are unchanged.

Define the entire even function
\[
 Q_\chi(w)=G_\chi(\tfrac12+w)G_{\bar\chi}(\tfrac12+w).
\]
Evenness follows from \eqref{eq:fe}. Its Taylor expansion therefore defines an entire $H_\chi$ by $Q_\chi(w)=H_\chi(w^2)$. Its order is at most $1/2$: for every $\eta>0$, the order bound on $Q_\chi$ gives $\log\max_{|z|=r}|H_\chi(z)|=O_\eta(r^{(1+\eta)/2})$. The critical-line theorem implies that every zero of $Q_\chi$ lies on the imaginary axis.

\begin{lemma}[Squared-coordinate representation]\label{lem:stieltjes}
Let $\rho=\tfrac12+i\gamma$ run over the nontrivial zeros of $L(s,\chi)$ with multiplicity, for primitive $\chi$ (including $q=1$). Then the zeros of $H_\chi$ are $-\gamma^2$, counted with exactly the multiplicity of this list. For $x>0$,
\begin{equation}\label{eq:F}
 F_\chi(x):=\frac{H_\chi'(x)}{H_\chi(x)}
 =\sum_{\rho}\frac{1}{x+\gamma^2}
 =\frac{1}{\sqrt{x}}\Rea\frac{G_\chi'}{G_\chi}\left(\tfrac12+\sqrt{x}\right).
\end{equation}
The positive series is absolutely convergent.
\end{lemma}
\begin{proof}
Conjugation exchanges the zeros of $G_\chi$ at $1/2+i\gamma$ with those of $G_{\bar\chi}$ at $1/2-i\gamma$. For $\gamma>0$, the multiplicity of $Q_\chi$ at $i\gamma$ is consequently the sum of the multiplicities of $L(s,\chi)$ at heights $\gamma$ and $-\gamma$. The local map $w\mapsto w^2$ is invertible there, giving that same multiplicity at $-\gamma^2$. At $0$, if the central zero of $G_\chi$ has order $r$, the product has order $2r$; the substitution gives order $r$ in $H_\chi$. This also explains why we use $\xi^2$, rather than $\xi$, for the trivial character. Since $H_\chi$ has order strictly less than one, Hadamard factorization is of genus zero with constant exponential factor. Thus, grouping positive squared ordinates $\lambda$ and their multiplicities $m_\lambda$,
\[
 H_\chi(z)=C z^r\prod_{\lambda>0}(1+z/\lambda)^{m_\lambda},
 \qquad \sum_{\lambda>0}m_\lambda/\lambda<\infty.
\]
Logarithmic differentiation, valid by local normal convergence away from the zeros, gives the middle series in \eqref{eq:F}. Differentiating $Q_\chi(w)=H_\chi(w^2)$ for positive real $w$ and using conjugation gives the last equality.
\end{proof}

\section{An exact rational-filter count certificate}
The main theorem is stated for an abstract positive discrete measure, so that the only number-theoretic dependence is Lemma~\ref{lem:stieltjes}.

Let $(t_j)_{j\in J}$ be a finite or countable list in $(0,1]$, repeated according to multiplicity, such that $M=\sum_jt_j<\infty$. For $0<\tau<1$ let
\[
 N_>(\tau)=\#\{j:t_j>\tau\},\qquad N_\ge(\tau)=\#\{j:t_j\ge\tau\}.
\]
Both are finite. Fix $n\ge1$ and distinct positive numbers $d_1,\ldots,d_n$. Define
\[
 D(t)=\prod_{i=1}^n(1+d_it),\quad d=D(1),\quad
 V_i=\sum_j\frac{t_j}{1+d_i t_j}.
\]
For $1\le k\le n$ define the binomial tail polynomial, with its usual extension $B_{n,0}=1$,
\begin{equation}\label{eq:binom}
 B_{n,k}(t)=\sum_{r=k}^n \binom{n}{r}t^r(1-t)^{n-r}.
\end{equation}
Its values at negative arguments are understood polynomially, not as probabilities. Set
\begin{align}
 c_i^{(k)}&=\frac{B_{n,k}(-1/d_i)}{(-1/d_i)\prod_{h\ne i}(1-d_h/d_i)},\label{eq:coeff}\\
 S_k&=\sum_{i=1}^n c_i^{(k)}V_i.\label{eq:Sk}
\end{align}

\begin{theorem}[Finite-sample certificate]\label{thm:count}
For any integers $1\le k_-\le k_+\le n$, set
\begin{equation}\label{eq:LU}
 \begin{split}
 L&=S_{k_+}-\frac{n}{k_+}B_{n-1,k_+-1}(\tau)M,\\
 U&=\frac{d}{B_{n,k_-}(\tau)}S_{k_-}.
 \end{split}
\end{equation}
Then
\[
 \boxed{\quad L\le N_>(\tau)\le N_\ge(\tau)\le U.\quad}
\]
In particular, if $\max(0,\lceil L\rceil)=\lfloor U\rfloor$, this common integer is the exact count, and there is no atom at $\tau$.
\end{theorem}
\begin{proof}
First observe the exact partial fraction identity
\begin{equation}\label{eq:identity}
 \frac{B_{n,k}(t)}{D(t)}=\sum_i c_i^{(k)}\frac{t}{1+d_it}.
\end{equation}
Multiplication by $D(t)$ yields two polynomials of degree at most $n$, both vanishing at $0$. At each $-1/d_i$ the right-hand polynomial has only its $i$th term nonzero, and \eqref{eq:coeff} makes it equal to the left. Their difference vanishes at $n+1$ distinct points and is identically zero. No derivatives or limiting node coalescence are used.

For $0\le t\le1$, $B_{n,k}(t)=\Pr(X\ge k)$ for $X\sim\mathrm{Bin}(n,t)$, so it is increasing in $t$, lies between $0$ and $1$, and obeys $B_{n,k}(t)\le nt/k$ by Markov's inequality. In fact the truncated first-moment identity yields the sharper
\begin{equation}\label{eq:moment}
 kB_{n,k}(t)\le \mathbb E[X\one_{X\ge k}]
 =nt B_{n-1,k-1}(t).
\end{equation}
The equality follows by the relation $r\binom{n}{r}=n\binom{n-1}{r-1}$. All sums of $B_{n,k}(t_j)/D(t_j)$ therefore converge absolutely. Summing \eqref{eq:identity} gives $S_k=\sum_j B_{n,k}(t_j)/D(t_j)$.

For $t\le\tau$, use \eqref{eq:moment} and $D(t)\ge1$ to see that
\[
 \frac{B_{n,k_+}(t)}{D(t)}-\frac n{k_+}B_{n-1,k_+-1}(\tau)t\le0.
\]
For $t>\tau$ the same expression is at most $1$. Hence it is a pointwise minorant of $\one_{t>\tau}$. For the majorant, monotonicity and $D(t)\le d$ imply
\[
 \frac{d B_{n,k_-}(t)}{D(t)B_{n,k_-}(\tau)}\ge\one_{t\ge\tau}.
\]
The denominator is positive since $\tau>0$ and $k_-\le n$. Summing both comparisons proves the result. The $t$ factor in the minorant correction is essential: an unweighted constant error over infinitely many atoms would diverge.
\end{proof}

\begin{corollary}[Dirichlet zero counts from the prime region]\label{cor:grh}
Let $a>1/4$, $T>0$, and define
\[
 t_\rho=\frac{a}{a+\gamma^2},\qquad \tau=\frac{a}{a+T^2},\qquad
 M=aF_\chi(a),\quad V_i=aF_\chi(a(1+d_i)).
\]
Under the proved critical-line theorem, the bounds \eqref{eq:LU} enclose, respectively, the strict and inclusive counts of strip zeros of $L(s,\chi)$ with $|\Im\rho|<T$ and $|\Im\rho|\le T$. These use only $n+1$ real-axis logarithmic-derivative samples at
\[
 s_0=\tfrac12+\sqrt a,\quad s_i=\tfrac12+\sqrt{a(1+d_i)}>1.
\]
They count multiplicities, including central zeros, and apply to imprimitive characters by using the inducing primitive character.
\end{corollary}
\begin{proof}
Apply Lemma~\ref{lem:stieltjes}, using $t_\rho/(1+d_i t_\rho)=a/[a(1+d_i)+\gamma^2]$. The inequalities $t_\rho>\tau$ and $t_\rho\ge\tau$ are equivalent to the stated strict and inclusive height conditions. The values $M,V_i$ are finite by the lemma.
\end{proof}

\section{Convergence, rates, and effective termination}
\begin{theorem}[Arbitrarily local sampling suffices]\label{thm:convergence}
Fix $a>1/4$, $0<\tau<1$ and any prescribed real sampling interval $(s_0,s_0+\ell)$, where $s_0=1/2+\sqrt a$ and $\ell>0$. One can choose a sequence of certificates with all $s_i$ in this interval such that
\[
 L_n\longrightarrow N_>(\tau),\qquad U_n\longrightarrow N_\ge(\tau).
\]
If there is no zero at the boundary $|\gamma|=T$, a finite certificate, evaluated with sufficiently accurate interval enclosures, determines the exact count. Effective termination assumes effectively given input parameters, for example rational $a,T$.
\end{theorem}
\begin{proof}
Choose a rational $c>0$ so small that $s_0+\ell>1/2+\sqrt{a(1+c)}$, and let $d_i=c i/n^3$, $1\le i\le n$. Then $0<d_i\le c$ and $1\le D_n(t)\le D_n(1)\le\exp(\sum_i d_i)\to1$ uniformly on $[0,1]$. Set $\delta_n=n^{-1/4}$ and $k_\pm=\lceil n(\tau\pm\delta_n)\rceil$ for all sufficiently large $n$ so that $1\le k_-\le k_+\le n$.

For each fixed $t$, binomial concentration gives $B_{n,k_+}(t)\to\one_{t>\tau}$ and $B_{n,k_-}(t)\to\one_{t\ge\tau}$. At $t=\tau$ this follows, for example, from Chebyshev's inequality, since $n\delta_n^2\to\infty$. The same argument at $n-1$ with threshold $k_+-1$ shows $B_{n-1,k_+-1}(\tau)\to0$. Also $B_{n,k_-}(\tau)\to1$. For both choices $k_\pm\ge n\tau/2$ eventually, hence $B_{n,k_\pm}(t)\le2t/\tau$. Since $\sum_jt_j=M<\infty$, dominated convergence applies to the countable sums, and \eqref{eq:LU} has the asserted limits.

When the two counts are equal to an integer $N$, the limiting endpoints approach $N$ from the correct sides. Eventually the lower endpoint exceeds $N-1$ and the upper endpoint is less than $N+1$. Any certified input approximations whose resulting propagated interval error tends to zero preserve this property eventually. The ceiling/floor test of Theorem~\ref{thm:count} terminates. This conclusion is not a test that is guaranteed to terminate at an actual boundary zero.
\end{proof}

\begin{proposition}[Gap-dependent error]\label{prop:rate}
Suppose no atom lies in $(\tau-g,\tau+g)$, where $0<g<\min(\tau,1-\tau)$. Take $k_\pm=\lceil n(\tau\pm g/2)\rceil$, $n-1\ge4/g$, and put $E=\exp(-(n-1)g^2/8)$. Write $N=N_>=N_\ge$. Then
\begin{align*}
 0\le N-L&\le\frac M\tau\left(1-\frac{1-E}{d}+E\right),\\
 0\le U-N&\le\frac M\tau\left(\frac{d}{1-E}-1+\frac{2dE}{1-E}\right).
\end{align*}
In particular, if the node spread is chosen so that $d-1=O(E)$, the count error is $O((M/\tau)E)$ for large $n$.
\end{proposition}
\begin{proof}
The ceiling rounding and $n-1\ge4/g$ imply
\[
 \frac{k_+-1}{n-1}\ge\tau+g/4,\quad
 \frac{k_--1}{n-1}\ge\tau-3g/4,
 \quad \frac{k_+}{n}\le\tau+3g/4,\quad
 \frac{k_-}{n}\le\tau-g/4.
\]
Hoeffding's inequality \cite{hoeffding} bounds any one-sided deviation of at least $g/4$ for the mean of $n$ or $n-1$ independent Bernoulli variables by $E$. In particular $B_{n-1,k_+-1}(\tau)\le E$ and $B_{n,k_-}(\tau)\ge1-E$. Each atom $t\ge\tau+g$ contributes at least $(1-E)/d$ to $S_{k_+}$; the other terms are nonnegative. For each $t\le\tau-g$, \eqref{eq:moment} and the second displayed bound imply $B_{n,k_-}(t)\le(n/k_-)t E\le2t E/\tau$. Thus
\[
 S_{k_+}\ge N(1-E)/d,\qquad S_{k_-}\le N+2ME/\tau.
\]
Finally $n/k_+\le1/\tau$ and $N\le M/\tau$. Substitute in \eqref{eq:LU} to obtain both bounds. A sufficiently small positive spread of distinct nodes makes $d$ as close to $1$ as desired, e.g. $d_i=E i/n^3$ (or rational lower approximations to these values).
\end{proof}

\paragraph{Other consequences.} Subtracting two certified inclusive counts at zero-free heights gives the exact number of zeros in the intervening symmetric annulus. Dovetail the count algorithms at rational heights to isolate every distinct positive absolute ordinate in arbitrarily short rational intervals, with its total multiplicity over the two signs: choose zero-free endpoints by testing many candidate rational endpoints in parallel, using the fact that the set of zero heights is discrete. More precisely, for any bounded height range with zero-free endpoints, for each mesh size enumerate rational partitions of that interval of maximum mesh at most the target size; run all their boundary certificates in parallel until a partition has certified counts. Such a partition exists by discreteness and density. The counts of its cells recover the zero counting measure to the given resolution (individual heights eventually separate). The construction does not distinguish positive from negative ordinates of a complex character, nor does it imply simplicity. A finite family of characters can share the same sampling nodes; alternatively sum their spectral transforms to count the combined multiset. A central zero is included automatically for $T>0$; by letting $T\downarrow0$ the count tends to the central multiplicity, but an unconditional finite stopping rule for determining that multiplicity does not follow from this argument alone.

\section{Certified inputs from the Euler product}
For real $s>1$, with $\Lambda(n)$ the von Mangoldt function, the classical identity is
\begin{equation}\label{eq:euler}
 \Rea\frac{G_\chi'(s)}{G_\chi(s)}=
 \frac12\log(q/\pi)+\frac12\psi\!\left(\frac{s+\kappa}{2}\right)
 -\sum_{m\ge2}\frac{\Lambda(m)\Rea\chi(m)}{m^s}
 +\one_{q=1}\left(\frac1s+\frac1{s-1}\right),
\end{equation}
where $\psi=\Gamma'/\Gamma$, $q=1$ means the primitive trivial character, and $\kappa=0$ in that case. It follows by logarithmic differentiation of the Euler product \cite{dlmf,chirre}. The series is absolutely convergent. For integer $X\ge3$ its discarded absolute tail is at most
\[
 \sum_{m>X}\frac{\log m}{m^s}\le
 X^{1-s}\left(\frac{\log X}{s-1}+\frac{1}{(s-1)^2}\right),
\]
by monotonicity of $u^{-s}\log u$ for $u\ge3$ and integration. The right side decreases to zero. All character values are zero or roots of unity and can be approximated rigorously from their finite character description. The gamma term is effectively computable, for example by the recurrence and controlled asymptotic expansion below. Thus when $a,T,d_i$ are rational (or effectively given with interval control), every sample can be enclosed by a rational interval of arbitrarily small width. No off-axis complex $L$ evaluations are needed. The real prime series is sufficient but not always the fastest numerical implementation.

For rational $\tau,d_i$, all coefficients in \eqref{eq:LU} are rational. Given interval enclosures $M\in[M_-,M_+]$ and $V_i\in[v_i^-,v_i^+]$, define the lower and upper bounds of a linear combination by choosing endpoints according to its coefficient's sign. The certified lower bound uses the lower linear sum for $S_{k_+}$ and $M_+$; the upper bound uses the upper linear sum for $S_{k_-}$. For example, if the width of each sample interval is at most $\varepsilon$, the additional enclosure width from this propagation is at most
\[
 \varepsilon\left(\sum_i|c_i^{(k_+)}|+\frac n{k_+}B_{n-1,k_+-1}(\tau)
 +\frac d{B_{n,k_-}(\tau)}\sum_i|c_i^{(k_-)}|\right).
\]
This is computable before sampling. Arbitrarily small node spread is an information theorem, not a numerical stability claim; the coefficients can become very large. For exact termination use, for example, a precision target making the displayed error at most $1/n$. No polynomial complexity bound is asserted.

\section{Reproducible numerical certification for zeta}\label{sec:numeric}
The archive contains a standard-library Python implementation (\texttt{certified\_zeta.py}) with outward fixed-point intervals and exact rational propagation. Its target is the known count at $T=8$, not a new zero-free region. Choose $a=64$, $n=16$, $\tau=1/2$, $k_-=7$, $k_+=10$ and $d_i=i/(4n^3)$. The base sample is $s_0=8.5$ and the remaining 16 lie just above it. The exact program output, rounded outwards, is
\[
 \boxed{-1.314658329293\le N_\zeta(8)\le0.186602311764.}
\]
Both strict and inclusive counts are therefore zero. The maximum sample interval width is $25750\cdot10^{-100}$. The serialized \texttt{zeta\_certificate.json} records all enclosures as scaled integer endpoints and the rational parameters. The verifier recomputes the final bracket; the analytic evaluator independently regenerates the input enclosures. We give the analytic tail bounds here so that the demonstration does not depend on an unverified zero table or an ad hoc floating-point tolerance.

\subsection{Fast real-axis zeta evaluation}
For $s>1$ let $\eta(s)=(1-2^{1-s})\zeta(s)$. For $m\ge0$ put
\[
 A_m(s)=\sum_{k=0}^m(-1)^k\binom{m}{k}(k+1)^{-s}
 =\frac1{\Gamma(s)}\int_0^\infty u^{s-1}e^{-u}(1-e^{-u})^m\,\dd u.
\]
The integral is obtained by expanding $(1-e^{-u})^m$ and evaluating gamma integrals. It shows $0\le A_m\le1$. Summing the geometric series inside the integral gives $\eta(s)=\sum_{m\ge0}2^{-m-1}A_m(s)$: indeed its kernel is $e^{-u}/(1+e^{-u})$, which expands to the alternating Dirichlet series for $s>1$. The first $J$ terms therefore have error at most $2^{-J}$. Differentiating the integral at $s>1$ is justified by dominated convergence on compact subintervals. If $Y$ has gamma density $u^{s-1}e^{-u}/\Gamma(s)$, then
\[
 |A_m'(s)|\le\mathbb E|\log Y|+|\psi(s)|
 \le2\bigl(s+1/(s-1)\bigr),
\]
using $|\log u|\le u+1/u$, $\psi(s)=\mathbb E\log Y$, and the gamma moments. Thus the omitted derivative has magnitude at most $2(s+1/(s-1))2^{-J}$. The finite double sum is collected into exact dyadic coefficients before evaluation, with $J=380$. The logarithmic derivative is
\[
 \frac{\zeta'(s)}{\zeta(s)}=\frac{\eta'(s)}{\eta(s)}-
 \frac{(\log2)2^{1-s}}{1-2^{1-s}}.
\]
Only real elementary functions are used.

\subsection{Digamma and rounding}
For positive $z$, Binet's integral formula differentiated once yields
\[
 \psi(z)=\log z-\frac1{2z}-2\int_0^\infty\frac{u}{(z^2+u^2)(e^{2\pi u}-1)}\,\dd u.
\]
Expand $1/(z^2+u^2)$ by the finite geometric identity. Integrating the first $p$ terms gives
\[
 \psi(z)=\log z-\frac1{2z}-\sum_{r=1}^p\frac{B_{2r}}{2r z^{2r}}+R_p,
 \qquad |R_p|\le\frac{|B_{2p+2}|}{(2p+2)z^{2p+2}}.
\]
For completeness, the integrated coefficient follows from expanding $(e^{2\pi u}-1)^{-1}$ into exponentials and using $2(2r-1)!\zeta(2r)/(2\pi)^{2r}=|B_{2r}|/(2r)$; the sign alternates by the sign of $B_{2r}$. The integral remainder is bounded by replacing its extra factor $1/(1+u^2/z^2)$ by $1$, which gives the first omitted magnitude. The code uses $z=x+100$ and $p=50$, then subtracts $\sum_{j=0}^{99}1/(x+j)$ by the digamma recurrence. Bernoulli numbers are computed by their rational recurrence. Machin's formula for $\pi$ uses alternating arctangent series with the first omitted term bounding the error.

Fixed-point intervals have endpoints in $10^{-100}\Z$. Addition, multiplication and division are outward-rounded by exact integer operations. Square root, logarithm and exponential use Python Decimal's correctly rounded operations at 112 significant digits; expanding each rounded endpoint by one neighboring representable decimal in the appropriate direction and rounding outwards to the fixed-point grid encloses the exact value. The only software trust beyond integer arithmetic and the proof is the implementation of these documented elementary operations \cite{decimal}. This is a reproducible interval computation, not a formally kernel-verified numerical proof. The separate \texttt{zero\_count.py} runs exact rational synthetic tests (including atoms at the threshold and repeated atoms); its optional mpmath demonstration is explicitly noncertified.

\begin{center}
\begin{tabular}{rrrrl}
\toprule
$n$ & strict count & inclusive count & & count enclosure (rounded display)\\
\midrule
32 & 2 & 2 && $[1.356633,2.374871]$\\
64 & 7 & 7 && $[6.177049,7.219724]$\\
32 & 2 & 5 && $[1.397932,5.380918]$\\
\bottomrule
\end{tabular}
\end{center}
The last row intentionally has a threshold atom of multiplicity three; the enclosure correctly retains the possible jump. Exact comparisons, not the displayed decimal approximations, are tested by the script.

\section{Novelty, limitations, and formalization route}
The classical parts are the functional equation, Hadamard product, Euler product, binomial concentration, and the principle of using a positive spectral measure. Standard zero counting uses the argument principle (e.g. Riemann--von Mangoldt and Turing methods); classical real-zero extraction also uses Euler--Rayleigh moment inequalities. Polynomial and rational spectral filtering for finite Hermitian eigenvalue counts is established, including work of Di Napoli, Polizzi and Saad \cite{dinapoli}. The independently derived contribution here is the specific two-sided rational-binomial certificate \eqref{eq:LU} on a possibly infinite counting measure with finite first moment, its derivative-free exact representation \eqref{eq:identity}, the strict/inclusive boundary limits, and the application to all Dirichlet $L$-functions via paired completions. The weighted correction is what makes the finite-count filter summable on the infinite zero spectrum. A limited web search for real-axis zero counting, Stieltjes transforms, Euler--Rayleigh bounds, and rational spectral filters did not locate this exact combination; this is not proof of bibliographic novelty or a priority claim. The illustrative zero count is not itself new. The note does not claim to solve the simple-zero conjecture or to outperform classical zero-search algorithms.

A Lean formalization can be split cleanly: (i) finite binomial polynomial identity, verified by degree and roots, (ii) finite Bernoulli probability inequalities and extension to summable series, (iii) the integer rounding rule, (iv) effective approximations, and (v) the analytic bridge from completed $L$-functions to a genus-zero product. The existing repository supplies critical-line nonvanishing but does not, by itself, certify this note's new lemmas. The included \texttt{FORMALIZATION.md} specifies the mathematical interfaces and dependencies; no new Lean build or proof term is claimed.

\begin{thebibliography}{99}\small
\bibitem{repo} L. Sernau et al., \emph{GRH for Dirichlet L-functions, in Lean 4}, 25 November 2026, \url{https://github.com/lsernau/grh-lean}. See \texttt{grh/GUIDE.md} and the readable edition source files cited in Section 1.
\bibitem{rubinstein} M. Rubinstein, \emph{Alp\"oge/Claude proof of the Generalized Riemann Hypothesis}, verification report and human-readable working guide, \url{https://www.math.uwaterloo.ca/~mrubinst/alpoge_claude_grh/}.
\bibitem{wiki} Wikipedia, \emph{Riemann hypothesis}, accessed 26 November 2026, \url{https://en.wikipedia.org/wiki/Riemann_hypothesis}.
\bibitem{davenport} H. Davenport, \emph{Multiplicative Number Theory}, 3rd ed., revised by H. L. Montgomery, Springer, 2000, chapters on Dirichlet functional equations and zeros.
\bibitem{dlmf} NIST Digital Library of Mathematical Functions, \S\S 25.2, 25.4, 25.10, 25.15 (zeta and Dirichlet $L$-functions), and \S\S 5.9, 5.11 (gamma and digamma), \url{https://dlmf.nist.gov/25.15} and \url{https://dlmf.nist.gov/5.11}.
\bibitem{chirre} A. Chirre and F. Gon\c{c}alves, \emph{Bounding the log-derivative of the zeta-function}, Math. Z. 300 (2022), 1041--1053, \href{https://doi.org/10.1007/s00209-021-02820-9}{doi:10.1007/s00209-021-02820-9}.
\bibitem{hoeffding} W. Hoeffding, \emph{Probability inequalities for sums of bounded random variables}, J. Amer. Statist. Assoc. 58 (1963), 13--30.
\bibitem{dinapoli} E. Di Napoli, E. Polizzi and Y. Saad, \emph{Efficient estimation of eigenvalue counts in an interval}, Numer. Linear Algebra Appl. 23 (2016), 674--692, \href{https://doi.org/10.1002/nla.2048}{doi:10.1002/nla.2048}.
\bibitem{decimal} Python Standard Library, \emph{decimal --- Decimal fixed-point and floating-point arithmetic}, \url{https://docs.python.org/3/library/decimal.html}.
\end{thebibliography}
\end{document}
